%% notes-tex/025-graph-reg-sweep/025-graph-reg-sweep.tex
%% [[experiments.025-solid-growth.graph-reg-sweep]]
%% https://github.com/Mjvolk3/torchcell/tree/main/notes-tex/025-graph-reg-sweep/025-graph-reg-sweep.tex
%%
%% Typeset counterpart of notes/experiments.025-solid-growth.graph-reg-sweep.md.
%%
%% Build:  make refresh  -> pull W&B, refill tables/, redraw the figures, convert, build
%%         make plots    -> figures/graph_reg_sweep_{ladder,curves,control}.pdf from the SVGs
%%         make          -> 025-graph-reg-sweep.pdf   (draft: status chips)
%%         make check    -> formatting / provenance / style gate
%%
%% Every number is read from W&B by experiments/025-solid-growth/scripts/graph_reg_sweep_readout.py
%% (results/graph_reg_sweep_runs.csv, graph_reg_sweep_history.csv, graph_reg_sweep_summary.json);
%% the tables under tables/ are written by the same script and never edited by hand; the
%% figures by graph_reg_sweep_plots.py from the same CSVs.

\documentclass[10pt,a4paper]{article}
\usepackage[draft]{tcdoc}
\usepackage{float}
\usepackage{amsmath}
\usepackage{amssymb}

\emergencystretch=3.5em

\title{\vspace{-1.5em}Soft Prior or Hard Mask\\
  {\large How the nine gene-gene graphs enter the cell graph transformer:
   the graph-regularization sweep on the 025 triples, read at three seeds per arm}}
\author{Michael Volk}
\date{\today}

\begin{document}
\maketitle
\thispagestyle{empty}

\begin{abstract}
\noindent
Nine gene-gene graphs enter layer 1 of the cell graph transformer by one of two
mechanisms: a soft prior, a divergence penalty of weight $\lambda$ that pulls each
head's attention toward its graph's normalized adjacency, or a hard mask that restricts
each head's attention to its graph's edges. The sweep runs one protocol (the 376{,}732
triples of experiment 010 on their pinned random split, a learnable gene table, AdamW at
a constant $2.5 \times 10^{-4}$, 30 epochs, four A40s per run) at $\lambda \in \{0,
10^{-5}, 10^{-4}, 10^{-3}, 10^{-2}, 10^{-1}, 1\}$, under the mask, and under the prior at
$\lambda = 10^{-3}$ toward degree-matched random rewirings of the same graphs, three
seeds per arm. Read at epoch 29, the mask sits with no penalty, $0.421 \pm 0.003$ against
$0.425 \pm 0.000$ held-out Pearson, and the prior rises with $\lambda$ without a peak
inside the ladder, to $0.447 \pm 0.001$ at $\lambda = 1$, above no penalty on every seed
by 0.021 to 0.023. Read at each run's best epoch, all nine arms lie within 0.010 of one
another (0.444 to 0.454): the prior moves the peak by at most 0.006 and moves when it
occurs, from epoch 15 or 16 without a penalty to epochs 23 to 26 at $\lambda = 1$, so most
of what it buys at the fixed reading is resistance to the decline the unpenalized model
shows after epoch 16. The attention contracts onto the graph at every $\lambda$ tested,
with edge recall at degree 0.78 at $10^{-5}$ and 0.915 from $10^{-2}$ up against 0.006
with no penalty, and saturates two decades below the $\lambda$ at which accuracy is
highest, so the contraction itself is not what carries the accuracy. The random-graph
control at $\lambda = 10^{-3}$, two complete seeds, sits at the no-penalty level,
$0.425 \pm 0.011$, where the biological graphs give $0.432 \pm 0.005$; at that $\lambda$
the effect is too small for two seeds to separate, and the control has not been run at
the $\lambda$ where the effect is largest. The hypothesis under test, that the graphs
help only when the model is optimized through the divergence and not when they are
imposed as a constraint, holds on the fixed-epoch reading and remains untested against
the alternative that any graph of the same degree sequence would do the same.
\end{abstract}

\vspace{0.5em}
\tccontents
\clearpage

%% notes-tex/025-graph-reg-sweep/sections/1-question.tex

\section{The question\secstatus{todo}}
\label{sec:question}

The cell graph transformer of experiment 010 attends over all 6{,}607 genes in every
layer. Nine gene-gene graphs, physical and regulatory interactions, TFLink, and the six
STRING channels (neighborhood, fusion, co-occurrence, co-expression, experimental,
database), are given to nine heads of layer 1, one graph per head. Two mechanisms can
hand a graph to a head. The \term{soft prior} adds to the training loss, for each head
$k$ with row-normalized adjacency $\tilde{A}_k$ and attention $\alpha_k$ on the
wild-type graph,
\begin{equation}
  \lambda \sum_{k=1}^{9} \frac{N}{|E_k|}
  \sum_{i} D_{\mathrm{KL}}\!\left(\tilde{A}_{k,i\cdot} \,\|\, \alpha_{k,i\cdot}\right),
  \label{eq:kl}
\end{equation}
with $N$ the number of genes and $|E_k|$ the edge count of graph $k$, so a head is pulled
toward spreading each gene's attention uniformly over that gene's neighbors and is free
to put mass anywhere. The \term{hard mask} sets the attention logits of every
non-edge (self-loops kept) to $-\infty$ before the softmax in layer 1, so a head can only
attend along the graph and needs no penalty and no $\lambda$. The two are the ends of one
axis: the mask is the $\lambda \to \infty$ limit of Eq.~\ref{eq:kl} on the support, but
not in the weights, since the mask lets a head weight its neighbors as it likes while the
prior pulls toward the uniform row of $\tilde{A}_k$.

Earlier readings put the two mechanisms far apart. The 010 model at $\lambda = 10^{-3}$
under a cosine schedule reached 0.456 to 0.464 held-out Pearson over three seeds, and on
the 025 build the same configuration replicated at 0.446 (GilaHyper job 1598, W\&B
\texttt{0yw7moue}), while the two hard-mask arms run then collapsed: 0.307 at epoch 1
and then near 0 (\texttt{7f1yrsq9}), or near 0 for 17 epochs (\texttt{4qmgkcgn}). Those
mask runs were under the cosine schedule with a different normalizer, so the collapse was
not a controlled comparison. The hypothesis this sweep was designed to test is stated in
the form it was proposed: \emph{the graphs help, but only when the model is optimized
through the divergence; imposing them as a mask does not help}. A second question rides
on the first: if the prior helps, is it the biology of the nine graphs or would any graph
with the same degree sequence condition training the same way?

\notesec{Design}
Every arm is the constant-rate control \texttt{cgt\_s0\_r\_kl\_ctrl\_013} with one key
changed (Table~\ref{tab:design}). The control is the 376{,}732 triples of experiment 010
(subset S0 of the 025 build) on 010's pinned random split, 301{,}386 training and
37{,}673 validation triples, a learnable 1{,}000-dimensional gene table, the label
normalizer fit on the training split, AdamW at $2.5 \times 10^{-4}$ with no schedule,
batch 256 over four GPUs, 30 epochs. The runs are Delta \texttt{gpuA40x4} jobs, one node
and four A40s per run, 24 h clocks, the 554 GB LMDB staged onto node-local NVMe at the
start of each job (\file{experiments/025-solid-growth/scripts/delta_cgt.slurm},
\file{experiments/025-solid-growth/scripts/delta_submit_sweep.sh}; jobs 22055147 to
22055169, 22056267 to 22056269, 22080055, 22080057 and 22107791). The prior arms peak at
6.05 GB of allocated GPU memory per card, the mask at 2.76 GB, because only the
penalized layer materializes its attention.

%% SOURCE: hand-authored from experiments/025-solid-growth/scripts/delta_submit_sweep.sh and the
%% configs cgt_s0_r_kl_ctrl_013, cgt_s0_r_mask_028, cgt_s0_r_kl_rand_031; seed counts from
%% experiments/025-solid-growth/results/graph_reg_sweep_summary.json (graph_reg_sweep_readout.py).
\begin{table}[htbp]
\centering
\footnotesize
\caption{The arms. Every arm is \texttt{cgt\_s0\_r\_kl\_ctrl\_013} with the one change
named; the control itself is the ladder's $\lambda = 10^{-3}$ point. Complete runs are
those with 30 logged epochs on W\&B at the pull of 2026-09-27; a seed listed as partial
was still running.}
\label{tab:design}
\begin{tabular}{llll}
\toprule
arm & change on the control & graph target & complete seeds \\
\midrule
no penalty & \texttt{graph\_reg\_lambda=0} & none & 1, 2, 3 \\
prior ladder & \texttt{graph\_reg\_lambda} $\in \{10^{-5}, 10^{-4}, 10^{-2}, 10^{-1}, 1\}$ & nine biological graphs, layer 1 & 1, 2, 3 each; $10^{-5}$: 1, 3 \\
control & none & nine biological graphs, layer 1, $\lambda = 10^{-3}$ & 1, 2, 3 \\
hard mask & \texttt{cgt\_s0\_r\_mask\_028}: mask on, $\lambda = 0$ & nine biological graphs, layer 1 & 1, 2, 3 \\
random graphs & \texttt{cgt\_s0\_r\_kl\_rand\_031}: rewiring on & nine rewired graphs, layer 1, $\lambda = 10^{-3}$ & 1, 2; 3 partial (epoch 16) \\
\bottomrule
\end{tabular}
\end{table}

\notesec{Readings}
Every run is read twice. The \term{fixed-epoch reading} is validation Pearson at epoch
29, the protocol's last epoch, the same optimization state for every arm. The
\term{max-over-epochs reading} is the largest validation Pearson over the 30 epochs, an
upward-biased order statistic reported with the epoch at which it occurred; it is the
number a best-checkpoint deployment would see. Arm values are the mean and sample
standard deviation over seeds; a paired comparison is by seed, since the split is pinned
and the seed sets the initialization and the batch order for every arm alike.

\notesec{Vocabulary}
Terms used in the figures and tables, in alphabetical order.

%% SOURCE: hand-authored; definitions follow torchcell/models/equivariant_cell_graph_transformer.py
%% (compute_graph_regularization_loss) and torchcell/trainers/int_transformer_cell.py (gradient probe,
%% edge recovery).
\begin{table}[htbp]
\centering
\footnotesize
\caption{Controlled vocabulary.}
\label{tab:vocab}
\begin{tabular}{p{3.2cm}p{13.6cm}}
\toprule
term & meaning \\
\midrule
divergence & Eq.~\ref{eq:kl} without its $\lambda$: the logged penalty divided by $\lambda$, summed over the nine heads and normalized by edge count, computed on the wild-type graph. Comparable across the ladder; not logged at $\lambda = 0$ or under the mask, where the model skips the computation. Reported from the validation pass at epoch 29. \\
edge recall at degree & For each gene $i$ of graph $k$ with degree $d_i$, the fraction of $i$'s neighbors among the $d_i$ largest entries of $\alpha_{k,i\cdot}$; averaged over genes, then over the nine heads. It is 1 by construction under the mask and near $1/N$ for attention unrelated to the graph. Logged at epochs 9, 19 and 29. \\
gradient probe & On the first batch of epochs 0, 1, 2, 5, 10 and 20, one extra backward pass per loss term; the $\ell_2$ norm of the parameter gradient of the point loss and of the penalty, before clipping, and their ratio. \\
point loss & Mean squared error of the prediction against the z-scored interaction label. \\
random graphs, degree-matched & Each of the nine graphs rewired by seeded double-edge swaps (five swaps per edge) that keep every gene's degree, in and out for the directed graphs; the rewiring seed is the run seed. The prior and the edge diagnostics are computed against the rewired graphs. \\
S0 & The 376{,}732 measured triples of experiment 010, as stored in the 025 build; the pool of every arm here. \\
\bottomrule
\end{tabular}
\end{table}

%% notes-tex/025-graph-reg-sweep/sections/2-ladder.tex

\section{The ladder: accuracy against the prior's weight\secstatus{todo}}
\label{sec:ladder}

\tcfig{figures/graph_reg_sweep_ladder.pdf}{The sweep at three seeds per arm, read
from W\&B by \file{experiments/025-solid-growth/scripts/graph_reg_sweep_readout.py}.
(a) Held-out Pearson on the 37{,}673 validation triples: filled points are the
fixed-epoch reading (epoch 29) of each seed with the arm mean as a bar and a line
through the prior ladder; hollow points are each seed's max over epochs. Gray is no
penalty, orange the prior toward the biological graphs (darker at higher $\lambda$), red
the hard mask, purple the prior toward degree-matched random graphs at $\lambda =
10^{-3}$. (b) Edge recall at degree at epoch 29, mean over the nine heads. (c) The
divergence of layer-1 attention from the nine graphs at epoch 29, validation pass, log
scale. (d) Gradient norm of the penalty over that of the point loss on the probe batch,
epoch 0 filled and epoch 20 hollow, log scale; the two arms without a penalty are at
0.}{fig:ladder}

\begin{table}[htbp]
\centering
\footnotesize
\caption{Arm means over seeds ($\pm$ sample sd) at the pull of 2026-09-27. Pearson is
held-out Pearson on the validation triples, at epoch 29 and as the max over epochs;
point loss and train Pearson are at epoch 29; divergence is the validation-pass value at
epoch 29; edge recall is at degree, mean over the nine heads, epoch 29. The highest
Pearson in each reading is bold. Generated by \texttt{graph\_reg\_sweep\_readout.py}.}
\label{tab:arms}
%% SOURCE: experiments/025-solid-growth/scripts/graph_reg_sweep_readout.py (W&B zhao-group/torchcell_025-solid-growth_equivariant_cell_graph_transformer, pulled 2026-10-11 01:26 UTC) -- GENERATED, do not edit

\begin{tabular}{lrrrr}
\toprule
arm & $n$ & Pearson, epoch 29 & Pearson, max over epochs & Pearson, at min val loss \\
\midrule
no penalty ($\lambda = 0$) & 3 & 0.425 $\pm$ 0.000 & 0.448 $\pm$ 0.002 & 0.440 $\pm$ 0.005 \\
KL $\lambda = 10^{-5}$ & 3 & 0.422 $\pm$ 0.005 & 0.445 $\pm$ 0.003 & 0.442 $\pm$ 0.003 \\
KL $\lambda = 10^{-4}$ & 3 & 0.419 $\pm$ 0.010 & 0.446 $\pm$ 0.003 & 0.442 $\pm$ 0.001 \\
KL $\lambda = 10^{-3}$ (control) & 3 & 0.432 $\pm$ 0.005 $\uparrow$ & 0.451 $\pm$ 0.004 $\uparrow$ & 0.444 $\pm$ 0.003 \\
KL $\lambda = 10^{-2}$ & 3 & 0.438 $\pm$ 0.002 $\uparrow$ & 0.453 $\pm$ 0.002 $\uparrow$ & 0.447 $\pm$ 0.005 \\
KL $\lambda = 10^{-1}$ & 3 & 0.445 $\pm$ 0.009 $\uparrow$ & \textbf{0.454 $\pm$ 0.006} & \textbf{0.453 $\pm$ 0.006} $\uparrow$ \\
KL $\lambda = 1$ & 3 & \textbf{0.447 $\pm$ 0.001} $\uparrow$ & 0.452 $\pm$ 0.003 $\uparrow$ & 0.443 $\pm$ 0.009 \\
hard mask, layer 1 & 3 & 0.421 $\pm$ 0.003 $\downarrow$ & 0.446 $\pm$ 0.003 & 0.442 $\pm$ 0.003 \\
KL $\lambda = 10^{-3}$, random graphs & 3 & 0.420 $\pm$ 0.012 & 0.447 $\pm$ 0.002 & 0.444 $\pm$ 0.006 \\
KL $\lambda = 10$ & 3 & 0.444 $\pm$ 0.008 $\uparrow$ & 0.452 $\pm$ 0.004 $\uparrow$ & 0.442 $\pm$ 0.005 \\
KL $\lambda = 100$ & 3 & 0.443 $\pm$ 0.004 $\uparrow$ & 0.451 $\pm$ 0.009 & 0.440 $\pm$ 0.004 \\
KL $\lambda = 10^{-1}$, random graphs & 5 & 0.430 $\pm$ 0.006 & 0.437 $\pm$ 0.003 $\downarrow$ & 0.434 $\pm$ 0.004 \\
KL $\lambda = 1$, random graphs & 5 & 0.427 $\pm$ 0.004 $\uparrow$ & 0.435 $\pm$ 0.004 $\downarrow$ & 0.430 $\pm$ 0.006 $\downarrow$ \\
hard mask, directed & 3 & 0.416 $\pm$ 0.012 $\downarrow$ & 0.448 $\pm$ 0.001 & 0.443 $\pm$ 0.002 \\
hard mask, 2 hops, directed & 3 & 0.419 $\pm$ 0.006 & 0.442 $\pm$ 0.003 $\downarrow$ & 0.440 $\pm$ 0.001 \\
hard mask, 3 hops, directed & 3 & 0.417 $\pm$ 0.003 $\downarrow$ & 0.445 $\pm$ 0.003 $\downarrow$ & 0.439 $\pm$ 0.004 $\downarrow$ \\
KL $\lambda = 1$, 2-hop target & 3 & 0.438 $\pm$ 0.006 $\uparrow$ & 0.442 $\pm$ 0.004 $\downarrow$ & 0.437 $\pm$ 0.004 $\downarrow$ \\
KL $\lambda = 1$, 3-hop target & 3 & 0.429 $\pm$ 0.012 & 0.442 $\pm$ 0.006 & 0.434 $\pm$ 0.004 $\downarrow$ \\
KL $\lambda = 1$, symmetric target & 3 & 0.439 $\pm$ 0.010 $\uparrow$ & 0.449 $\pm$ 0.003 & 0.446 $\pm$ 0.007 \\
hard mask, layers 1--2 & 3 & 0.419 $\pm$ 0.004 $\downarrow$ & 0.446 $\pm$ 0.001 $\downarrow$ & 0.440 $\pm$ 0.004 \\
hard mask, layers 3--4 & 3 & 0.415 $\pm$ 0.007 $\downarrow$ & 0.445 $\pm$ 0.001 $\downarrow$ & 0.442 $\pm$ 0.002 \\
hard mask, layers 1--4 & 3 & 0.426 $\pm$ 0.009 & 0.446 $\pm$ 0.002 & 0.445 $\pm$ 0.003 \\
KL $\lambda = 1$, layers 1--2 & 3 & 0.437 $\pm$ 0.009 $\uparrow$ & 0.446 $\pm$ 0.006 & 0.440 $\pm$ 0.007 \\
KL $\lambda = 1$, layers 3--4 & 3 & 0.445 $\pm$ 0.007 $\uparrow$ & 0.449 $\pm$ 0.004 & 0.441 $\pm$ 0.011 \\
no penalty, 60 epochs & 3 & 0.415 $\pm$ 0.003 & 0.443 $\pm$ 0.004 & 0.440 $\pm$ 0.001 \\
hard mask, 60 epochs & 3 & 0.419 $\pm$ 0.009 & 0.446 $\pm$ 0.001 $\uparrow$ & 0.440 $\pm$ 0.004 \\
KL $\lambda = 1$, 60 epochs & 3 & 0.441 $\pm$ 0.005 $\uparrow$ & 0.450 $\pm$ 0.006 & 0.445 $\pm$ 0.011 \\
no penalty, composite embedding & 3 & -0.006 $\pm$ 0.004 & 0.073 $\pm$ 0.061 & -0.002 $\pm$ 0.001 \\
KL $\lambda = 1$, composite embedding & 3 & 0.421 $\pm$ 0.007 $\uparrow$ & 0.423 $\pm$ 0.007 $\uparrow$ & 0.423 $\pm$ 0.007 $\uparrow$ \\
\bottomrule
\end{tabular}

\end{table}

\notesec{The mask is no penalty}
On the fixed-epoch reading the hard mask gives $0.421 \pm 0.003$ and no penalty
$0.425 \pm 0.000$ (Table~\ref{tab:arms}); by seed the mask is below no penalty on all
three, by 0.001, 0.007 and 0.003. On the max-over-epochs reading the two are
$0.446 \pm 0.003$ and $0.448 \pm 0.002$, the mask below on two seeds of three. The mask
arm also fits the training triples least, 0.612 against 0.624 train Pearson at epoch
29. Restricting the support of layer-1 attention to the nine graphs therefore changes
nothing a validation triple can see. The collapse of the two earlier mask runs did not
recur under this protocol: all three seeds train normally
(Figure~\ref{fig:curves}a).

\notesec{The prior rises with $\lambda$ and does not peak inside the ladder}
At epoch 29 the ladder reads 0.420, 0.419, 0.432, 0.438, 0.445 and 0.447 at
$\lambda = 10^{-5}$ through 1 (Table~\ref{tab:arms}). The two smallest weights sit
at or below no penalty, both seeds of $10^{-5}$ below all three no-penalty seeds; from
$10^{-3}$ up every seed is above its no-penalty partner, by 0.005 to 0.013 at $10^{-3}$
and by 0.021, 0.023 and 0.021 at $\lambda = 1$. The mock-up drawn before the runs
(\file{notes/assets/drawio/FigS-graph-regularization-sweep.gen.py}) placed a peak
between the two extremes; the measurement places the maximum at the top of the ladder,
with $10^{-1}$ and 1 within one standard deviation of each other ($0.445 \pm 0.009$,
$0.447 \pm 0.001$). Whether the curve turns over at $\lambda = 10$ or 100 is not
measured.

\notesec{The max-over-epochs reading compresses the ladder}
Read at each run's best epoch, the nine arms span 0.444 to 0.454, and the best arm
($\lambda = 10^{-1}$, $0.454 \pm 0.006$) is 0.006 above no penalty ($0.448 \pm 0.002$).
At $\lambda = 1$ the paired differences against no penalty shrink from 0.021 to 0.023 at
epoch 29 to 0.007, 0.003 and 0.003 at the best epoch, still positive on all three seeds.
What changes with $\lambda$ is less the height of the peak than its timing: the
no-penalty seeds peak at epochs 15, 16 and 16, the mask at 15, 17 and 18, $\lambda = 1$
at 23, 24 and 26 (Table~\ref{tab:runs}). The fixed-epoch gain of the prior is therefore
mostly the unpenalized model's decline from epoch 16 to epoch 29 (0.441 to 0.425 in the
seed mean) that the prior at high $\lambda$ does not show (0.435 at 15, 0.447 at 29).

%% notes-tex/025-graph-reg-sweep/sections/3-training.tex

\section{What the prior does to training\secstatus{todo}}
\label{sec:training}

\notesec{Support contraction happens at every $\lambda$, and saturates early}
Without a penalty, layer-1 attention has nothing to do with the graphs: edge recall at
degree is 0.006 at every diagnostic epoch, the value of attention unrelated to the
adjacency (Figure~\ref{fig:sweep}c). The smallest weight tested already moves it to
0.782 at epoch 29 (0.759 at epoch 9), $10^{-4}$ to 0.866, $10^{-3}$ to 0.906, and from
$10^{-2}$ up it sits at 0.914 to 0.915 with no further change to $\lambda = 1$. The
divergence tells the same story with more resolution (Figure~\ref{fig:sweep}d): from
1{,}540 at $10^{-5}$ through 571 and 337 to about 300 at $10^{-2}$, $10^{-1}$ and 1
(300, 297, 298), a floor the penalty cannot push below at any weight tested. Both
diagnostics are flat from $10^{-2}$ to 1 while the fixed-epoch accuracy still rises from
0.438 to 0.447. How close the attention sits to the graph is set by $\lambda = 10^{-2}$;
what the two larger weights add is not more contraction.

\notesec{The gradient budget}
The gradient probe (Figure~\ref{fig:sweep}e) measures the norm of each loss term's
gradient on one batch before clipping. The point-loss gradient at epoch 0 is a property
of the seed, not the arm: 7.4, 32.3 and 2.8 for seeds 1, 2 and 3 in every arm that
shares the initialization (the mask arm, with a different layer-1 kernel, reads 5.7,
32.9 and 2.7). The penalty's gradient scales with $\lambda$ as it must, 1.1 at
$10^{-3}$, 10.9 at $10^{-2}$, 109 at $10^{-1}$ and 1{,}097 at 1 in the seed mean at
epoch 0, and grows over training while the point-loss gradient settles to 1.0 to 4.3:
at epoch 20 the ratio is 1.1 to 1.5 at $10^{-3}$, 15 to 24 at $10^{-2}$, 86 to 146 at
$10^{-1}$ and 680 to 1{,}550 at $\lambda = 1$. The two best arms of the ladder train
under a total gradient that is almost entirely the penalty. The trainer clips the total
gradient norm at 10, so at $\lambda \geq 10^{-1}$ every update is rescaled by two to
three orders of magnitude before the optimizer sees it.

Hypothesis (untested): AdamW normalizes each parameter's update by its running gradient
scale, so a uniform rescaling of the whole gradient leaves the update nearly unchanged,
and the penalty's gradient is concentrated in the parameters that shape layer-1
attention (the layer-1 query and key projections and, through them, the gene table).
Under that reading the arms at $\lambda \geq 10^{-1}$ train the rest of the network on
the point loss much as the no-penalty arm does, which is consistent with their training
Pearson matching it (0.625 and 0.624 against 0.624), while layer 1 is held on the graph
for the whole run. The mechanism is not measured here; a per-parameter-group norm on
the probe batch would measure it.

\notesec{The prior is not a fit-reducing regularizer at the weights that help}
Training Pearson at epoch 29 (Figure~\ref{fig:sweep}h) is 0.624 with no penalty, 0.636
at $10^{-3}$, 0.613 at $10^{-2}$, 0.625 at $10^{-1}$ and 0.624 at 1, so the arms that
generalize best at epoch 29 fit the training triples as well as the arm that generalizes
worst. The two smallest weights fit the training set better than any other arm, 0.657
at $10^{-5}$ and 0.663 at $10^{-4}$, and generalize no better than no penalty. The
validation point loss (Figure~\ref{fig:sweep}i) reaches its minimum near epoch 13 in
every arm (0.80 to 0.82) and then rises; at epoch 29 it is 0.867 with no penalty and
0.869 under the mask against 0.826 and 0.825 at $\lambda = 10^{-1}$ and 1. The prior at
high weight does not stop the model from fitting the training triples; it slows the
growth of held-out error after the point-loss minimum.

Hypothesis (untested): the late decline of the unpenalized model is layer-1 attention
drifting away from any fixed structure as the gene table specializes to the training
triples, and a prior strong enough to hold layer 1 in place removes that degree of
freedom without touching the rest. The mask removes the same freedom in a different
way, by fixing the support and leaving the weights free, and it declines exactly as no
penalty does, so if this hypothesis is right the part of layer 1 that drifts is the
weighting within the support, not the support.

%% notes-tex/025-graph-reg-sweep/sections/4-control.tex

\section{The random-graph control: biology or any target\secstatus{todo}}
\label{sec:control}

\tcfig{figures/graph_reg_sweep_control.pdf}{The prior at $\lambda = 10^{-3}$ toward
the nine biological graphs (orange, three seeds), toward degree-matched random
rewirings of the same graphs (purple, two complete seeds; the cross is seed 3's max over
epochs at epoch 16, partial), and no penalty (gray, three seeds). (a) Fixed-epoch
reading with the arm mean as a bar. (b) Validation Pearson by epoch, mean over seeds
with a $\pm$ sd band.}{fig:control}

The control arm keeps everything of the $\lambda = 10^{-3}$ control except the graphs:
each of the nine is rewired before the model sees it, by double-edge swaps that keep
every gene's degree, five swaps per edge, seeded by the run seed
(\texttt{cgt\_s0\_r\_kl\_rand\_031}, \texttt{torchcell.graph.rewire}). The prior pulls
the heads toward the rewired graphs and the edge diagnostics are scored against them.
The arm was written to run at the best $\lambda$ of the sweep; the sweep's best was not
known at submission, so it ran at $10^{-3}$, the 010 operating point.

\notesec{The heads learn the random graphs, less completely}
Edge recall at degree against the rewired graphs reaches 0.814 at epoch 29 and the
divergence 983 (Figure~\ref{fig:ladder}b, c), against 0.906 and 337 for the biological
graphs at the same $\lambda$. A random graph of the same degree sequence is a target the
attention can and does move toward, but it does not get as close in 30 epochs. Whether
that is because the biological graphs are more compatible with the point loss, or
because the rewired graphs are less clustered and so harder to represent with a
low-rank attention, is not separated by anything logged.

\notesec{At $\lambda = 10^{-3}$ the accuracy question is underpowered}
On the fixed-epoch reading the two complete random-graph seeds give 0.417 and 0.433,
against 0.429 and 0.438 for the biological graphs on the same seeds and 0.424 and 0.425
with no penalty (Figure~\ref{fig:control}a). The biological graphs are above the random
ones on both seeds, by 0.012 and 0.005; the random ones straddle no penalty. At the
max-over-epochs reading the three arms are 0.451, 0.447 and 0.448. The effect being
tested at this $\lambda$ is the 0.005 to 0.013 that the biological prior adds over no
penalty at epoch 29, and two seeds of a control whose own seed spread is 0.011 cannot
resolve it. The question the control was designed to answer, whether the biology or the
conditioning carries the gain, is open, and it is answerable only where the gain is
large: at $\lambda = 1$ the biological prior is 0.021 to 0.023 above no penalty on every
seed, so three seeds of the rewired graphs at $\lambda = 1$ would separate a 0.02 effect
from the 0.001 to 0.011 seed spread seen here. That arm is one override away
(\texttt{model.graph\_regularization.graph\_reg\_lambda=1} on
\texttt{cgt\_s0\_r\_kl\_rand\_031}) and has not been submitted.

%% notes-tex/025-graph-reg-sweep/sections/5-reading.tex

\section{Reading, and what is still missing\secstatus{todo}}
\label{sec:reading}

\notesec{On the hypothesis}
Stated as proposed: the graphs help, but only when the model is optimized through the
divergence; imposing them as a mask does not help. The mask half holds on both readings
and all three seeds: the mask is indistinguishable from no penalty (0.421 against 0.425
at epoch 29, 0.446 against 0.448 at the best epoch), and it trains without the collapse
the earlier cosine-schedule runs showed. The prior half holds on the fixed-epoch reading
from $\lambda = 10^{-3}$ up, with the gain growing to 0.021 to 0.023 at $\lambda = 1$ and
no turnover inside the ladder, and holds weakly on the best-epoch reading (0.003 to
0.007 at $\lambda = 1$, positive on every seed). The word ``help'' needs its reading
attached: a best-checkpoint deployment gains at most 0.006 from the prior; a fixed-budget
run gains 0.02, because the prior removes the decline after epoch 16 rather than raising
the peak. What the sweep does not establish is that the help comes from the biology of
the graphs. The random-graph control ran only at $\lambda = 10^{-3}$, where the effect is
within two seeds' noise, so ``networks help'' is measured as ``a strong prior toward these
graphs helps'' and the alternative, that a strong prior toward any graph of the same
degree sequence helps as much, is untested.

\notesec{What the mock-up predicted and what was measured}
The figure design (\file{notes/assets/drawio/FigS-graph-regularization-sweep.gen.py})
drew six panels. Panel a, contraction with $\lambda$: measured, and it saturates at
$10^{-2}$. Panel b, accuracy peaking between the extremes: measured, and there is no
interior peak; the maximum is at the top of the ladder. Panel c, the penalty carrying the
gradient before the point loss can: measured, and the penalty dominates the gradient at
every probe epoch for $\lambda \geq 10^{-2}$, not only early. Panels d and e, discovery
of edges outside a head's prior and overruling of cataloged edges by the digenic
interaction, need the layer-1 attention of the saved checkpoints, which are on Delta
\texttt{/scratch} and have not been read; the edge-recovery precision logged at $k =
32$ (0.21 for every prior arm from $10^{-3}$ up, 0.006 with no penalty) is the nearest
logged quantity and is not the discovery statistic. Panel f, the random-graph control:
measured at the wrong $\lambda$.

\notesec{Open items}
\begin{enumerate}
\item The random-graph arm at $\lambda = 1$, three seeds. The one run that decides
  between conditioning and biology.
\item $\lambda = 10$ and 100, three seeds each, to find the turnover the ladder did
  not reach. Memory is unchanged (the penalty is computed at every $\lambda > 0$).
\item Seed 2 of $\lambda = 10^{-5}$ (Delta job 22055161) has no run on W\&B; the job's
  log on Delta says why. The arm is read at two seeds.
\item Seed 3 of the random-graph arm (job 22056269) was at epoch 16 at the pull; the
  tables refresh with \texttt{make refresh}.
\item Panels d and e from the checkpoints: pull the best-Pearson checkpoint of one seed
  per arm from Delta and score layer-1 attention against TFLink-only and
  STRING-experimental-only edges and against Costanzo digenic $\varepsilon$.
\item A per-parameter-group gradient probe, to test the hypothesis of
  Section~\ref{sec:training} about where the penalty's gradient lands.
\end{enumerate}

\begin{table}[htbp]
\centering
\footnotesize
\caption{Every run read, one row per seed. Epochs is the number of validation epochs
logged; Pearson at epoch 29 and the max over epochs with its epoch; the gradient-probe
ratio (penalty over point loss) at epoch 0; the W\&B run id in project \texttt{torchcell\_025-solid-growth\_equivariant\_cell\_graph\_transformer}.
Generated by \texttt{graph\_reg\_sweep\_readout.py}.}
\label{tab:runs}
%% SOURCE: experiments/025-solid-growth/scripts/graph_reg_sweep_readout.py (W&B zhao-group/torchcell_025-solid-growth_equivariant_cell_graph_transformer, pulled 2026-10-01 21:40 UTC) -- GENERATED, do not edit

\begin{tabular}{llrrrrl}
\toprule
arm & seed & epochs & Pearson, ep 29 & max (epoch) & gradient ratio, ep 0 & W\&B run \\
\midrule
no penalty ($\lambda = 0$) & 1 & 30 & 0.4243 & 0.4476 (16) & 0.00 & \texttt{zvtj3waw} \\
no penalty ($\lambda = 0$) & 2 & 30 & 0.4248 & 0.4460 (15) & 0.00 & \texttt{4vgeyw2i} \\
no penalty ($\lambda = 0$) & 3 & 30 & 0.4253 & 0.4493 (16) & 0.00 & \texttt{mmiydyws} \\
KL $\lambda = 10^{-5}$ & 1 & 30 & 0.4213 & 0.4456 (13) & 0.00 & \texttt{rldsrcnr} \\
KL $\lambda = 10^{-5}$ & 3 & 30 & 0.4180 & 0.4425 (14) & 0.00 & \texttt{7fx9njte} \\
KL $\lambda = 10^{-4}$ & 1 & 30 & 0.4189 & 0.4473 (16) & 0.01 & \texttt{r2i3i9xq} \\
KL $\lambda = 10^{-4}$ & 2 & 30 & 0.4297 & 0.4486 (15) & 0.00 & \texttt{h8zhfq2d} \\
KL $\lambda = 10^{-4}$ & 3 & 30 & 0.4096 & 0.4431 (6) & 0.04 & \texttt{oxh1ug5l} \\
KL $\lambda = 10^{-3}$ (control) & 1 & 30 & 0.4292 & 0.4498 (16) & 0.14 & \texttt{6dvbzy8w} \\
KL $\lambda = 10^{-3}$ (control) & 2 & 30 & 0.4378 & 0.4471 (12) & 0.04 & \texttt{ppm9kl4w} \\
KL $\lambda = 10^{-3}$ (control) & 3 & 30 & 0.4304 & 0.4548 (14) & 0.40 & \texttt{22u65gax} \\
KL $\lambda = 10^{-2}$ & 1 & 30 & 0.4358 & 0.4524 (27) & 1.45 & \texttt{u1hdm0y0} \\
KL $\lambda = 10^{-2}$ & 2 & 30 & 0.4386 & 0.4510 (19) & 0.37 & \texttt{eufg8mlm} \\
KL $\lambda = 10^{-2}$ & 3 & 30 & 0.4386 & 0.4541 (14) & 4.01 & \texttt{0vgat9et} \\
KL $\lambda = 10^{-1}$ & 1 & 30 & 0.4528 & 0.4590 (23) & 14.47 & \texttt{lz1rd90p} \\
KL $\lambda = 10^{-1}$ & 2 & 30 & 0.4473 & 0.4542 (23) & 3.69 & \texttt{df8jt886} \\
KL $\lambda = 10^{-1}$ & 3 & 30 & 0.4349 & 0.4478 (13) & 40.20 & \texttt{5pz6vc5y} \\
KL $\lambda = 1$ & 1 & 30 & 0.4456 & 0.4545 (26) & 146.51 & \texttt{367oc8pv} \\
KL $\lambda = 1$ & 2 & 30 & 0.4480 & 0.4487 (23) & 37.17 & \texttt{ds928dmp} \\
KL $\lambda = 1$ & 3 & 30 & 0.4461 & 0.4521 (24) & 399.95 & \texttt{zxme1g1p} \\
hard mask, layer 1 & 1 & 30 & 0.4232 & 0.4430 (18) & 0.00 & \texttt{g35yf77n} \\
hard mask, layer 1 & 2 & 30 & 0.4181 & 0.4491 (15) & 0.00 & \texttt{1pv2lpf9} \\
hard mask, layer 1 & 3 & 30 & 0.4222 & 0.4458 (17) & 0.00 & \texttt{tpixauk7} \\
KL $\lambda = 10^{-3}$, random graphs & 1 & 30 & 0.4169 & 0.4459 (13) & 0.15 & \texttt{usgha5w6} \\
KL $\lambda = 10^{-3}$, random graphs & 2 & 30 & 0.4329 & 0.4500 (19) & 0.04 & \texttt{9tqgl66k} \\
KL $\lambda = 10^{-3}$, random graphs & 3 & 30 & 0.4104 & 0.4456 (14) & 0.41 & \texttt{zuw5z7k4} \\
KL $\lambda = 10^{-1}$, random graphs & 1 & 30 & 0.4343 & 0.4343 (29) & 14.78 & \texttt{u6he7prk} \\
KL $\lambda = 10^{-1}$, random graphs & 2 & 30 & 0.4300 & 0.4376 (12) & 3.71 & \texttt{h7wumv3p} \\
KL $\lambda = 10^{-1}$, random graphs & 3 & 30 & 0.4203 & 0.4341 (19) & 40.87 & \texttt{w7wt6bmo} \\
KL $\lambda = 10^{-1}$, random graphs & 4 & 30 & 0.4294 & 0.4412 (21) & 2.14 & \texttt{7e9caxfo} \\
KL $\lambda = 10^{-1}$, random graphs & 5 & 30 & 0.4372 & 0.4374 (24) & 22.20 & \texttt{jpmuaaah} \\
KL $\lambda = 1$, random graphs & 1 & 30 & 0.4321 & 0.4366 (28) & 149.41 & \texttt{l6hv7cub} \\
KL $\lambda = 1$, random graphs & 2 & 30 & 0.4298 & 0.4342 (14) & 37.41 & \texttt{awnp9vvt} \\
KL $\lambda = 1$, random graphs & 3 & 30 & 0.4260 & 0.4404 (22) & 405.90 & \texttt{gtmbi3fw} \\
KL $\lambda = 1$, random graphs & 4 & 30 & 0.4256 & 0.4348 (28) & 20.97 & \texttt{urwi2lxu} \\
KL $\lambda = 1$, random graphs & 5 & 30 & 0.4212 & 0.4287 (19) & 223.91 & \texttt{n0vst095} \\
\bottomrule
\end{tabular}

\end{table}


\end{document}
